\subsection{\label{sec:AnisotropicSymanzik}Anisotropic Symanzik gauge action}

\subsubsection{\label{sec:AnisoTreeLevel}The tree-level action}

For anisotropic lattices with bare gauge anisotropy $\xi_0$, the tree-level Symanzik gauge action is
\begin{equation}
\begin{split}
&S_G=\beta \sum _{n}\left[\sum _{i<j}\frac{1}{\xi _0}\left(c_P P_{ij}(n)+c_R R_{ij}(n)\right)+\sum _{i}\xi _0\left(c_P P_{i4}(n)+c_R R_{i4}(n)\right)\right]\\
&c_P=1,\;\;c_R=-\frac{1}{20},\;\;\beta =\frac{10N}{3g^2}\quad\left(\beta=\frac{10}{g^2}\ \text{for SU(3)}\right)\\
\end{split}
\end{equation}
where $P_{\mu\nu}$ and $R_{\mu\nu}$ are the plaquette and the $2\times 1$ rectangle in the $\mu$-$\nu$ plane. Planes made of only spatial directions carry the weight $1/\xi _0$, and planes containing the temporal direction carry $\xi _0$. The field level kernels use the inverse convention, so $q=1/\xi _0$ is passed to the kernels.

\subsubsection{\label{sec:AnisoTwistedLoop}The twisted (parallelogram) loop term}

The one-loop improved action also contains the twisted loop with six links $[\mu ,\nu ,\rho ,-\mu ,-\nu ,-\rho ]$ (also called a parallelogram or "rectangle without right angles"). There are four direction triples: the purely spatial triple $(1,2,3)$, and the three triples $(1,2,4)$, $(1,3,4)$, $(2,3,4)$ containing the temporal direction.

A common mistake is to weight the whole loop by "whether the triple contains the temporal direction". In the weak field expansion of one loop,
\begin{equation}
\begin{split}
&1-\frac{1}{N}{\rm Re} {\rm tr} \left[U_{\mu\nu\rho}(n)\right]\sim F_{\mu\nu}^2+F_{\mu\rho}^2+F_{\nu\rho}^2+{\rm cross\;terms}\\
\end{split}
\end{equation}
and the cross terms cancel after summing the four sign variants, so every twisted loop contributes to all the three planes of its triple. Defining $Q_{123}$ as the sum of the purely spatial triple and $Q_{\sigma\sigma\tau}=Q_{124}+Q_{134}+Q_{234}$, their expansions are
\begin{equation}
\begin{split}
&Q_{123}\sim F_{12}^2+F_{13}^2+F_{23}^2\\
&Q_{\sigma\sigma\tau}\sim \left(F_{12}^2+F_{13}^2+F_{23}^2\right)+2\left(F_{14}^2+F_{24}^2+F_{34}^2\right)\\
\end{split}
\end{equation}
Therefore a time-containing twisted loop is not a purely temporal operator, it still contains one spatial $F_{ij}^2$. Requiring the spatial planes to carry $1/\xi _0$ and the temporal planes to carry $\xi _0$, the minimal generalization of the isotropic term $S_T^{iso}=\beta c_T(Q_{123}+Q_{\sigma\sigma\tau})$ with the correct continuum limit is
\begin{equation}
\begin{split}
&S_T^{ani}=\beta c_T\left[\left(\frac{2}{\xi _0}-\xi _0\right)Q_{123}+\xi _0 Q_{\sigma\sigma\tau}\right]\\
\end{split}
\end{equation}
which reduces to the isotropic one at $\xi _0=1$. Equivalently, with the separated operators
\begin{equation}
\begin{split}
&Q_s=Q_{123},\;\;Q_t=\frac{1}{2}\left(Q_{\sigma\sigma\tau}-Q_{123}\right),\;\;S_T^{ani}=2\beta c_T\left(\frac{1}{\xi _0}Q_s+\xi _0 Q_t\right)\\
\end{split}
\end{equation}
which shows why the subtraction of $Q_{123}$ is unavoidable. In the kernel convention $q=1/\xi _0$, the weights are
\begin{equation}
\begin{split}
&w_{123}=2q-\frac{1}{q},\;\;w_{\sigma\sigma\tau}=\frac{1}{q}\\
\end{split}
\end{equation}
for the purely spatial and the time-containing triples, and the energy and the force kernels use exactly the same weights.

\subsubsection{\label{sec:AnisoOneLoopCaveat}Limitations}

The twisted term above is a continuum-matched (heuristic) extension, not a complete one-loop anisotropic Symanzik improvement. A true one-loop action must split all operator classes,
\begin{equation}
\begin{split}
&S_g=\beta \sum _n\left[c_{P,s}P_{ij}+c_{P,t}P_{i4}+c_{R,s}R_{ij}+c_{R,st}^{(s)}R_{ij4}+c_{R,st}^{(t)}R_{i44}+c_{T,s}Q_s+c_{T,t}Q_t+\cdots \right]\\
\end{split}
\end{equation}
with coefficients $c_a=c_a^{(0)}+\alpha _s c_a^{(1)}(\xi _0,N_f)$ rematched by anisotropic lattice perturbation theory, and in general two tadpole factors $u_s$ and $u_t$. The isotropic MILC coefficients
\begin{equation}
\begin{split}
&c_R=-\frac{1}{20u_0^2}\left[1-(0.6264-1.1746N_f)\ln u_0\right],\;\;c_T=\frac{(0.0433-0.0156N_f)\ln u_0}{u_0^2}\\
\end{split}
\end{equation}
are strictly valid only at $\xi _0=1$. For the strict implementation, use the tree-level action with $c_T=0$ and $c_R=-1/20$.
